El sistema puede desplegarse en cuatro escenarios distintos según el contexto de uso:
desde un único ordenador para desarrollo y pruebas hasta un clúster Kubernetes para
producción en la nube. La tabla~\ref{tab:escenarios_despliegue} resume cada opción.

\begin{table}[H]
  \centering
  \caption{Escenarios de despliegue de Voctomix~2.0.}
  \label{tab:escenarios_despliegue}
  \small
  \begin{tabularx}{\textwidth}{|l|X|l|}
    \hline
    \textbf{Escenario} & \textbf{Requisito} & \textbf{Caso de uso} \\
    \hline
    Un PC (nativo)   & Linux, GStreamer y Python       & Pruebas y desarrollo \\
    Dos PCs (nativo) & Dos equipos Linux en red local  & Pruebas y desarrollo \\
    Docker Compose   & Docker instalado                & Entorno portable sin instalación \\
    Kubernetes       & Minikube o clúster externo      & Infraestructura orquestada \\
    \hline
  \end{tabularx}
\end{table}

Los dos primeros se denominan escenarios nativos porque ejecutan los procesos directamente
sobre el sistema operativo del equipo, sin virtualización. Los dos siguientes emplean
contenedores, lo que elimina la necesidad de instalar GStreamer, Python o sus dependencias
en la máquina anfitriona y garantiza un comportamiento idéntico en cualquier equipo.
La contenerización fue un objetivo prioritario del proyecto por dos razones complementarias.
Docker simplifica la adopción del sistema al reducir el despliegue a un único comando,
sin configuración manual del entorno. Kubernetes, por su parte, responde a una necesidad
concreta del contexto en el que se desarrolla este trabajo: los proyectos de investigación
CyberNEMO de la UPM utilizan Kubernetes como infraestructura de producción, por lo
que Voctomix~2.0 debe poder integrarse en ese entorno.

\section{Despliegue nativo: un PC y dos PCs}
\label{sec:despliegue_nativo}

El despliegue nativo ejecuta los procesos de voctocore y voctogui directamente sobre
Linux, sin ninguna capa intermedia. Es el escenario de referencia para el desarrollo
y las pruebas funcionales del sistema.

\begin{itemize}

  \item \textbf{Escenario de un PC.}
  Este fue el escenario utilizado durante el desarrollo del proyecto y el que permitió
  comprender en profundidad la arquitectura del sistema. El arranque sigue tres pasos:
  primero se lanza \texttt{voctocore.py}, que abre el puerto de control~(9999) y los
  puertos de entrada de señal (10000--10005, uno por fuente), y comienza a escuchar
  conexiones; a continuación se inyectan las señales de vídeo y audio en esos puertos
  mediante FFmpeg; por último se arranca \texttt{voctogui.py}, que se conecta a
  voctocore a través de \texttt{localhost} (configurado como \texttt{host=localhost}
  en el fichero de configuración de voctogui, al ser la dirección del propio equipo),
  y muestra en la interfaz todas las fuentes activas. Al cerrar la interfaz, los
  procesos auxiliares se terminan de forma limpia de manera automática.

  \item \textbf{Escenario de dos PCs.}
  En este escenario, la carga de trabajo se reparte entre dos equipos conectados por
  red local: el PC servidor ejecuta voctocore y todos los servicios de apoyo, mientras
  que el PC del realizador ejecuta únicamente voctogui. Para que voctogui pueda
  conectarse a voctocore, el fichero de configuración de voctogui debe indicar la
  dirección IP del servidor en el campo \texttt{host}, en lugar de \texttt{localhost}.
  El script de arranque del realizador acepta esta IP como variable de entorno
  (\texttt{IP\_SERVER}), lo que evita modificar el fichero de configuración en cada
  despliegue.

\end{itemize}
